\section{Discussion}
\label{sec:discussion}

The central finding --- that within-game z-score normalization reveals 4--5 genuinely
independent dimensions where raw PCA shows only one --- has a clear practical implication:
the choice of normalization method is not a technical detail but a theoretical commitment.
Raw PCA answers the question ``how much of the variation between games is explained by
overall quality?'' (answer: most of it). Within-game normalization answers the different
question ``how much of the variation in design character is explained by orthogonal
design dimensions?'' (answer: the five TIGER dimensions explain $\geq 80\%$).

\subsection{What Each Dimension Represents}

\textbf{Tension (T).} The axes that load on this dimension --- time pressure, catch-up
difficulty, penalty severity, and elimination risk --- collectively encode the psychological
pressure experienced by the player. This maps to Knizia's oft-cited design principle that
meaningful decisions require real stakes: without Tension, choices are preferences rather
than decisions. The empirical orthogonality of Tension from Interaction confirms that
pressure can be generated by the game system (solo, cooperative, or racing structures)
independently of direct player competition.

\textbf{Interaction (I).} Interaction loads on axes measuring direct conflict, negotiation
opportunity, player targeting, and theft/blocking. It is structurally independent of
Tension because competitive games can be either high-Interaction (direct combat) or
low-Interaction (parallel optimization). The separation of I and T is the empirical
confirmation of the design-community intuition that ``mean'' and ``tense'' are different
game properties.

\textbf{Game-Architecture (G).} This dimension captures production depth, engine
complexity, phase count, and rulebook length --- the infrastructure required to
support the game's decision space. High-G games (Lacerda's catalog) require significant
setup and explanation investment. The architectural depth is orthogonal to both Tension
and Elegance, confirming that complexity is not equivalent to either pressure or opacity.

\textbf{Elegance (E).} Elegance loads on rule parsimony, mechanical unity, and the
ratio of strategic depth to rulebook length. It corresponds directly to the
German-school design ideal: more decisions per rule, fewer exceptions, strong
thematic integration. The negative correlation between E and G (low but consistent)
reflects the design tension between depth-through-architecture and depth-through-elegance.

\textbf{Range (R).} Range captures the breadth of viable strategies, player-count
flexibility, and scaling across skill levels. It is the least correlated with overall
game weight, confirming that accessibility and replayability are determined by design
choices orthogonal to both complexity and elegance. High-R games (Dominion, 7~Wonders)
achieve breadth through procedural generation of the game state rather than through
simplified rules.

\subsection{Limitations}

The normalization method assumes that every game has meaningful variation across axes
(i.e., $\sigma_g > 0$ for all games). Games that score near-uniformly across all axes
--- perhaps because they are extremely generic designs --- have near-zero denominators
and are excluded from the normalized analysis. This affects approximately 2\% of the
corpus and does not materially affect the PCA results.
